\documentclass[11pt]{article}

\usepackage[margin=1in]{geometry}
\usepackage{amsmath,amssymb}
\usepackage{array}
\usepackage{url}

\newcolumntype{L}[1]{>{\raggedright\arraybackslash}p{#1}}
\urlstyle{tt}
\DeclareRobustCommand{\code}[1]{\ifmmode\text{\texttt{\detokenize{#1}}}\else\texttt{\detokenize{#1}}\fi}
\newcommand{\Dpres}{D_{\mathrm{pres}}}

\title{Dimensionality Extension Plan\\Belief-State (Bayesian) Adaptive Fractionation}
\author{Internal repository design note}
\date{\today}

\begin{document}
\maketitle

\section{Purpose}

This note defines the belief-state extension to the adaptive fractionation solver across two stages of increasing Bayesian fidelity. Both stages retain a 5-coordinate finite-horizon state.

The baseline (current master branch) state is:
\begin{equation}
\text{state}_{\mathrm{current}} = (t,\, o_t,\, D_t).
\end{equation}

\textbf{Stage A} (implemented) extends the state to include a running Gaussian belief:
\begin{equation}
\text{state}_{\mathrm{A}} = (t,\, o_t,\, D_t,\, \mu_t,\, \sigma_t),
\end{equation}
where $(\mu_t, \sigma_t)$ tracks a running mean and standard deviation of the patient's PTV--OAR overlap distribution, updated after each observed fraction via Welford's online algorithm. Branch transition probabilities are computed from a Normal predictive distribution $\mathcal{N}(\mu_t, \sigma_t^2)$.

\textbf{Stage B} (NIG extension) retains the same 5-coordinate state but replaces the Normal predictive with the theoretically correct posterior predictive derived from the Normal-Inverse-Gamma (NIG) conjugate model. The key observation enabling this is that the full NIG posterior is a deterministic function of the sufficient statistics $(n, \bar s, M_2)$, all of which are already encoded in the Stage~A state $(\mu_t, \sigma_t, t)$ together with the fixed population prior. No additional state dimensions are required: Stage~B is implemented by changing how branch-probability tables are computed, not by expanding the state.

\section{Scope and Non-Goals}

\subsection{In Scope}
\begin{itemize}
\item Replace the Normal branch-probability tables with NIG posterior-predictive (Student-$t$) tables, indexed by fraction number.
\item Use the prior-corrected posterior mean $\mu_n = (\kappa_0\mu_0 + n\bar s)/(\kappa_0+n)$ as the predictive centre rather than the raw sample mean.
\item Update planning-time initialisation to use the full NIG posterior predictive rather than a Normal approximation.
\item Estimate and introduce the two additional NIG prior hyperparameters $\mu_0$ (prior mean) and $\kappa_0$ (prior precision weight) from population-level estimation on the patient cohort, stored alongside the existing \code{DEFAULT_ALPHA} and \code{DEFAULT_BETA} constants.
\item Preserve Welford's belief-update rule unchanged: it is equivalent to updating the NIG sufficient statistics.
\item Update the planning-scan termination criterion in \code{precompute_plan} to use a fixed clinically safe overlap bound rather than a Gaussian percentile heuristic.
\end{itemize}

\subsection{What Remains Unchanged}
\begin{itemize}
\item The state coordinates $(t, o_t, D_t, \mu_t, \sigma_t)$ and their discretisation grids.
\item The belief-update operator: Welford's online algorithm.
\item Per-fraction action bounds $d_t \in [d_{\min}, d_{\max}]$ with the same action discretisation by \code{dose_steps}.
\item Feasibility screening: infeasible states are assigned a hard sentinel value using the same convention as Stage~A.
\item Terminal prescription logic: only $D_F = \Dpres$ is feasible at treatment end; terminal underdose and overdose remain hard-infeasible.
\end{itemize}

\subsection{Out of Scope}
\begin{itemize}
\item Changes to the clinical objective form beyond existing penalty logic.
\item UI redesign in \code{app.py} beyond minimal parameter exposure.
\end{itemize}

\section{Required Mathematical Changes}

\subsection{Stage A (Implemented): Bellman Recursion}

The Stage~A recursion is:
\begin{equation}
V_t^{\mathrm{A}}(o,D,\mu,\sigma) =
\max_d \Bigl[
-\hat c_t(o,d)
+ \mathbb{E}_{o' \sim \mathcal N(\mu,\sigma^2)}
V_{t+1}^{\mathrm{A}}\!\left(o', D+d,\; \mathcal U_{\mathrm{W}}(\mu,\sigma,o',t)\right)
\Bigr].
\end{equation}

The belief update $\mathcal U_{\mathrm{W}}$ is Welford's running-moment increment. Let $n_t = t$ be the number of observations accumulated at fraction $t$; then:
\begin{align}
\mu' &= \frac{n_t\,\mu + o'}{n_t + 1}, \\[4pt]
(\sigma')^2 &= \frac{n_t\,\sigma^2 + (o'-\mu)(o'-\mu')}{n_t + 1}.
\end{align}
This is the population-variance (ddof$=0$) form of Welford's increment, consistent with the codebase convention and with the denominator convention of the 2025 paper.

\noindent\textit{Initialisation at planning time.} At decision time for fraction $f$ (with $f$ real observations $o_1,\dots,o_f$): set $\mu_{\mathrm{start}} = \bar o_f$ (population mean of observations) and $\sigma_{\mathrm{start}}$ from the MAP estimator (equations~11--13 of the clinical trial paper) applied to all $f$ observations, using prior parameters $(\alpha_0, \beta_0)$. This gives a prior-regularised $\sigma$ that is nonzero even at $f=1$.

\noindent\textit{Branch representative.} For each overlap bin $j$ with edges $(l_j, u_j]$, use the grid centre \code{volume_space[j]} as the concrete scalar fed into $\mathcal{U}_{\mathrm{W}}$.

\subsection{Stage A: Grid Specification}

The $\mu$ grid uses five non-uniform segments over $[0, 30]$\,cc, with finest resolution in the typical-patient range and progressively coarser resolution toward extreme values:
\begin{center}
\small
\begin{tabular}{llll}
\hline
\textbf{Segment} & \textbf{Range (cc)} & \textbf{Points} & \textbf{Step (cc)} \\
\hline
1 & $[0.00,\ 1.00]$  & 38 & $\approx 0.027$ \\
2 & $[1.05,\ 4.00]$  & 39 & $\approx 0.078$ \\
3 & $[4.10,\ 10.00]$ & 42 & $\approx 0.144$ \\
4 & $[10.20,\ 16.00]$ & 15 & $\approx 0.414$ \\
5 & $[16.50,\ 30.00]$ & 16 & $\approx 0.900$ \\
\hline
\multicolumn{2}{l}{\textbf{Total}} & \textbf{150} & \\
\hline
\end{tabular}
\normalsize
\end{center}

The $\sigma$ grid uses 11 non-uniform points over $[0.05, 4.5]$\,cc, with fine resolution in the range where most patients fall and coarser resolution in the tail:
\begin{center}
\small
\begin{tabular}{llll}
\hline
\textbf{Segment} & \textbf{Range (cc)} & \textbf{Points} & \textbf{Step (cc)} \\
\hline
1 & $[0.05,\ 0.70]$ & 5 & $\approx 0.163$ \\
2 & $[0.80,\ 1.80]$ & 3 & $\approx 0.500$ \\
3 & $[2.00,\ 4.50]$ & 3 & $\approx 1.250$ \\
\hline
\multicolumn{2}{l}{\textbf{Total}} & \textbf{11} & \\
\hline
\end{tabular}
\normalsize
\end{center}
This gives $150 \times 11 = 1{,}650$ belief states. The upper limit of 4.5\,cc is set to avoid clipping the most variable patient in the 58-patient cohort (patient~3, MAP $\hat\sigma = 4.05$\,cc). The fixed overlap state space \code{_VOLUME_SPACE} covers $[0, 44]$\,cc at 441 uniformly spaced points (step 0.1\,cc, matching TPS volume discretisation). The 44\,cc bound is hardcoded and independent of $\sigma_{\max}$, so the bin width stays exactly 0.1\,cc regardless of future $\sigma$ grid changes; 44\,cc comfortably covers $\mu_{\max} + 3\sigma_{\max} = 30 + 3\times 4.5 = 43.5$\,cc and the tail probability beyond 44\,cc is negligible for any belief on the grid. Branch probabilities fold left- and right-tail mass into the boundary bins so that probabilities sum to~1 for every belief-grid point. This tail-folding argument for Stage~A relies on the Normal predictive being negligibly thin beyond 44\,cc. For Stage~B, before committing the 44\,cc upper bound, sweep all $(t, \mu_{\mathrm{grid}}, \sigma_{\mathrm{grid}})$ triples under the fitted NIG hyperparameters and record the Student-$t$ predictive tail mass above 44\,cc. If the maximum folded mass is non-negligible, enlarge \code{_VOLUME_SPACE} (e.g.\ to 50 or 60\,cc) and re-run the benchmark to measure policy sensitivity; alternatively, document the approximation error explicitly if the 44\,cc cap is operationally required.

\noindent\textit{Numerical invariant.} \code{_VOLUME_SPACE} is fixed and shared across all beliefs and all fractions. Using patient-specific or belief-specific overlap ranges is explicitly prohibited: when a belief's $\mu$ is far from the patient's current range, a narrow range assigns nearly all probability mass to a single edge bin, making the DP plan as if all future overlaps equal that edge regardless of the true belief. A fixed wide range avoids this.

\subsection{Stage B: NIG Posterior from Sufficient Statistics}

The NIG model places a Normal-Inverse-Gamma prior on the unknown overlap-distribution parameters $(\mu, \sigma^2)$:
\[
\mu,\,\sigma^2 \;\sim\; \mathrm{NIG}(\mu_0,\,\kappa_0,\,\alpha_0,\,\beta_0).
\]
Given $n$ i.i.d.\ observations with sample mean $\bar s$ and sum of squared deviations $M_2 = \sum_{i=1}^n(o_i - \bar s)^2$, the posterior is $\mathrm{NIG}(\mu_n, \kappa_n, \alpha_n, \beta_n)$ with the closed-form updates:
\begin{align}
\kappa_n &= \kappa_0 + n, \label{eq:kappa-update} \\[3pt]
\mu_n    &= \frac{\kappa_0\,\mu_0 + n\,\bar s}{\kappa_n}, \label{eq:mu-update} \\[3pt]
\alpha_n &= \alpha_0 + \tfrac{n}{2}, \label{eq:alpha-update} \\[3pt]
\beta_n  &= \beta_0 + \tfrac{M_2}{2} + \frac{\kappa_0\,n\,(\bar s - \mu_0)^2}{2\,\kappa_n}. \label{eq:beta-update}
\end{align}

\noindent\textit{Key reduction.} The Stage~A state encodes all four sufficient statistics:
\begin{align*}
n     &= t \quad\text{(absolute 1-based fraction index, tracked as \code{observation_count = int(fraction_number)})}, \\
\bar s &= \mu_t \quad\text{(Welford running mean, a state coordinate)}, \\
M_2   &= n\,\sigma_t^2 \quad\text{(recoverable from $\sigma_t$ and $t$)}.
\end{align*}
Therefore $(\kappa_n, \mu_n, \alpha_n, \beta_n)$ are computable deterministically at any DP node from $(\mu_t, \sigma_t, t)$ and the fixed prior constants $(\mu_0, \kappa_0, \alpha_0, \beta_0)$. No additional state coordinates are required.

\noindent Two practical caveats. First, the grid quantisation of $\sigma_t$ introduces the same discretisation error as in Stage~A and does not degrade Stage~B relative to Stage~A. Second, the Welford update in \code{_hypothetical_belief_grid_indices} and \code{_bellman_expectation_full_grid} applies a floor $\sigma' \ge \sigma_{\min} = 0.05$\,cc before snapping to the grid; the effective $M_2$ entering the $\beta_n$ computation is therefore $n \times \max(\sigma_t^2,\, \sigma_{\min}^2)$, not $n \times \sigma_t^2$. This is a minor effect for typical patients but must be matched exactly in the $\beta_n$ computation to keep Stage~B consistent with Stage~A's Welford dynamics.

\subsection{Stage B: Posterior Predictive Distribution}

Given the posterior $\mathrm{NIG}(\mu_n, \kappa_n, \alpha_n, \beta_n)$, the posterior predictive distribution for the next observation is:
\begin{equation}
\label{eq:predictive-student-t}
o_{t+1} \mid (\mu_t,\,\sigma_t,\,t) \;\sim\; \mathrm{Student\text{-}t}_{2\alpha_n}\!\bigl(\mu_n,\; s_n\bigr),
\end{equation}
where the predictive scale is
\begin{equation}
\label{eq:predictive-scale}
s_n = \sqrt{\frac{\beta_n\,(\kappa_n+1)}{\alpha_n\,\kappa_n}},
\end{equation}
and all quantities on the right-hand side are computed from $(\mu_t, \sigma_t, t, \mu_0, \kappa_0, \alpha_0, \beta_0)$ via equations~\eqref{eq:kappa-update}--\eqref{eq:beta-update}.

\noindent\textit{Why Student-$t$.} When both the mean and variance of the overlap distribution are uncertain and marginalised out, the predictive is Student-$t$ rather than Normal. This gives heavier tails, most pronounced at early fractions where $\alpha_n$ is small and the degrees of freedom $2\alpha_n = 2\alpha_0 + t$ are low (near-Cauchy at $t=1$). Here ``near-Cauchy'' means a very low-degree-of-freedom Student-$t$: not literally a Cauchy unless the degrees of freedom equal 1, but close in the practical sense that it is much less concentrated than a Normal and assigns substantially more probability to extreme overlaps far from its centre. As $t$ grows, $2\alpha_n \to \infty$ and the Student-$t$ converges to a Normal; Stage~B and Stage~A become equivalent in the large-$t$ limit.

\noindent\textit{Finite-variance constraint.} Enforce $\alpha_0^{\mathrm{NIG}} \geq 1$ in the empirical Bayes optimisation, which guarantees degrees of freedom $2\alpha_n \geq 2+t$ at every fraction and therefore a finite-variance predictive throughout treatment. Note that $\alpha_0 = 1$ is the minimum acceptable value but remains the boundary case: at fraction~1 the predictive is Student-$t_3$, which has finite variance but still places roughly five times more probability mass beyond $3\times\text{scale}$ than a Normal. The fraction-1 predictive should be plotted and inspected for both representative and extreme patients regardless of where the optimum falls. If the unconstrained optimisation consistently prefers $\alpha_0 < 1$, this is a model-misspecification signal; report both the unconstrained and constrained optima so the discrepancy can be investigated before finalising the model.

\noindent\textit{Prior-corrected centre.} The predictive centre $\mu_n$ is a precision-weighted average of the population prior mean $\mu_0$ and the patient's running mean $\mu_t$. At $t=1$ with $\kappa_0=1$ (for illustration; see Section~3.6 on why $\kappa_0=1$ is not the recommended choice), prior and observation receive equal weight. With a small fitted $\kappa_0$, the predictive centre converges to $\mu_t$ after even one or two observations. As $t\to\infty$, $\mu_n\to\mu_t$ regardless of $\kappa_0$. This shrinkage is appropriate when the patient's overlap history is short and the individual estimate is unreliable; the magnitude of the shrinkage is controlled by fitting $\kappa_0$ from the cohort data.

\noindent\textit{Connection to the branch update.} The Welford update equations for $(\mu, \sigma)$ are algebraically identical to the NIG sufficient-statistics update for $(\bar s, M_2)$: both compute the same new mean and the same new sum of squared deviations. The belief-update operator $\mathcal U_{\mathrm{W}}$ is therefore already the correct NIG update and requires no modification in Stage~B.

\subsection{Stage B: Bellman Recursion}

The Stage~B recursion has the same state and structure as Stage~A:
\begin{equation}
V_t^{\mathrm{B}}(o,D,\mu,\sigma) =
\max_d \Bigl[
-\hat c_t(o,d)
+ \mathbb{E}_{o' \sim \mathrm{St}_{2\alpha_n}(\mu_n,\, s_n)}
V_{t+1}^{\mathrm{B}}\!\left(o', D+d,\; \mathcal U_{\mathrm{W}}(\mu,\sigma,o',t)\right)
\Bigr],
\end{equation}
where $\alpha_n$, $\mu_n$, and $s_n$ are computed on-the-fly from $(\mu, \sigma, t)$ via equations~\eqref{eq:kappa-update}--\eqref{eq:predictive-scale}, and $\mathcal U_{\mathrm{W}}$ is unchanged from Stage~A.

\subsection{Stage B: Implementation Strategy}

Stage~B requires three targeted changes to Stage~A and no structural changes to the DP.

\paragraph{Change 1: Fraction-indexed branch-probability tables.}

In Stage~A, \code{_P_BELIEF[mi, si, j]} is precomputed once at module load using Normal CDFs and is the same for all fractions. In Stage~B, the branch probability depends on the fraction index $t$ because the Student-$t$ degrees of freedom $2\alpha_n(t)$ and predictive scale $s_n(\mu,\sigma,t)$ both depend on $t$. The precomputed table becomes:
\begin{equation*}
\text{\code{_P_BELIEF_NIG[t, mi, si, j]}}
\;=\;
P_{\mathrm{St}}\!\Bigl(\text{bin }j \;\Big|\;
\mu_n\!\left(\mu_{\mathrm{grid}}[m_i], t\right),\;
s_n\!\left(\mu_{\mathrm{grid}}[m_i], \sigma_{\mathrm{grid}}[s_i], t\right),\;
2\alpha_n(t)
\Bigr),
\end{equation*}
computed using Student-$t$ CDF differences with the same left/right-tail folding as Stage~A. The table shape is $(N_F, N_\mu, N_\sigma, N_o)$; for a 5-fraction treatment with the current grids this is $5\times150\times11\times441 \approx 3.6\,\text{M}$ entries ($\approx 29$\,MB), a negligible addition to module-load memory.

In the current codebase, \code{_P_BELIEF} is a module-level global accessed directly inside \code{_bellman_expectation_full_grid} (not passed as a parameter). \code{_P_BELIEF_NIG} must follow the same pattern: a module-level array of shape $(N_F, N_\mu, N_\sigma, N_o)$ precomputed at module load. The first axis is \textit{0-based}: axis index 0 corresponds to $t=1$ (fraction~1), axis index $t-1$ corresponds to fraction~$t$. In the function body, replace \code{_P_BELIEF[:, :, j]} with \code{_P_BELIEF_NIG[observation_count - 1, :, :, j]}. \code{_bellman_expectation} takes branch probabilities as an explicit argument and is unaffected by this change.

\noindent\textit{Variable treatment lengths.} The public solver API accepts arbitrary \code{number_of_fractions}; the table must not be hard-wired to a single treatment length. At module load, precompute the table for \code{DEFAULT_NUMBER_OF_FRACTIONS} and cache the result in a module-level variable. If a caller requests a different treatment length, recompute for that length and update the cache. This is simpler than a fully general lazy-cache keyed by an arbitrary parameter and covers all current code paths, including the existing 3-fraction and 5-fraction test cases. Add regression tests exercising both.

\paragraph{Change 2: Planning-time initialisation.}

At decision time for fraction $f$, \code{current_overlap_probs} is currently computed by \code{probdist} (\code{helper_functions.py}), which returns Normal CDF differences without folding tail mass into the boundary bins; the function explicitly notes that its output does not sum to~1 if the distribution extends outside \code{_VOLUME_SPACE}. Under Stage~B, this must be replaced with a Student-$t$ predictive using parameters $(\mu_n, s_n, 2\alpha_n)$ evaluated at $n=f$ from the $f$ actual observations. The replacement \emph{must} fold left and right tails into the boundary bins, matching the behaviour of \code{_P_BELIEF_NIG}, so that planning-time and forward-DP branch probabilities are computed on the same probabilistic basis.

\noindent\textit{No $\sigma$ point estimate required.} The Student-$t$ predictive parameters $(\mu_n, s_n, \nu_n = 2\alpha_n)$ follow directly from the NIG sufficient statistics $(n, \bar s, M_2)$ and the fitted hyperparameters via equations~\eqref{eq:kappa-update}--\eqref{eq:predictive-scale}; the posterior does not need to be reduced to a scalar $\hat\sigma$ before the predictive can be evaluated. \code{std_calc} must not be called in the Stage~B planning-time path. If a $\sigma$-grid index is needed (for example to look up an initial DP value), the NIG posterior mode of $\sigma$ when $\sigma^2 \sim \mathrm{InvGamma}(\alpha_n, \beta_n)$ is $\sqrt{\beta_n/(\alpha_n + \tfrac{1}{2})}$. Do not use $\sqrt{\beta_n/(\alpha_n - 1)}$, which equals $\sqrt{E[\sigma^2]}$ (the square root of the posterior mean of $\sigma^2$, not the mode); at small $\alpha_n$ this overestimates $\hat\sigma$ considerably.

\paragraph{Prior hyperparameters: refitting required.}

Stage~A uses $(\alpha_0, \beta_0)$ (\code{DEFAULT_ALPHA}, \code{DEFAULT_BETA}) to regularise the MAP estimate of $\sigma$ via a numerical grid search in \code{std_calc}. Critically, \code{std_calc} assumes a \textit{Gamma}$(\alpha, \beta)$ prior on $\sigma$ directly (density $\propto \sigma^{\alpha-1}\exp(-\sigma/\beta)$), whereas the NIG model assumes an \textit{Inverse-Gamma}$(\alpha, \beta)$ prior on $\sigma^2$ (density $\propto (\sigma^2)^{-\alpha-1}\exp(-\beta/\sigma^2)$). These are materially different: with the current values, the Gamma prior on $\sigma$ has mode $(\alpha_0-1)\beta_0 \approx 0.06$\,cc, while the Inverse-Gamma prior on $\sigma^2$ with the same $(\alpha_0, \beta_0)$ numbers implies a $\sigma$-mode of $\sqrt{\beta_0/(\alpha_0+1)} \approx 0.61$\,cc. Reusing \code{DEFAULT_ALPHA} and \code{DEFAULT_BETA} directly in equation~\eqref{eq:beta-update} would silently produce incorrect $\beta_n$ values.

Before Stage~B is implemented, new constants \code{NIG_ALPHA_0} and \code{NIG_BETA_0} must be estimated from the patient cohort under the Inverse-Gamma convention. The recommended procedure is empirical Bayes: treat each patient's overlap sequence as i.i.d.\ draws from a Normal with unknown $(\mu, \sigma^2)$, place the NIG prior on $(\mu, \sigma^2)$, and maximise the marginal log-likelihood $\sum_i \log p(o_{i,1},\dots,o_{i,f_i} \mid \mu_0, \kappa_0, \alpha_0, \beta_0)$ over the full 58-patient cohort. Since $\kappa_0$ and $\mu_0$ interact with $\alpha_0$ and $\beta_0$ through the marginal (Student-$t$) likelihood, all four parameters should be fitted jointly with $\kappa_0$ included as a free parameter (bounds $[10^{-3}, 5]$). Do not fix $\kappa_0 = 1$: the 58-patient cohort has a wide spread of overlap means (0.01--21.18\,cc), and $\kappa_0 = 1$ over-shrinks the predictive centre of extreme patients toward the population mean, which can represent a qualitative regression relative to Stage~A at early fractions. After fitting, check the Stage~B predictive centre at fraction~1 for the most extreme patients (patient~5, $\mu \approx 21$\,cc; patient~3, $\sigma \approx 4.05$\,cc) to confirm that shrinkage is clinically reasonable. Report the fitted $\kappa_0$ value prominently. The marginal likelihood for a single patient with observations $o_{1:n}$ and prior $({\mu_0, \kappa_0, \alpha_0, \beta_0})$ is a product of Student-$t$ densities with sequentially updated parameters, which follows directly from applying equations~\eqref{eq:kappa-update}--\eqref{eq:beta-update} at each step.

\code{DEFAULT_ALPHA} and \code{DEFAULT_BETA} must be retained unchanged for \code{std_calc} and Stage~A initialisation. The two additional constants $\mu_0$ and $\kappa_0$ are also required and should be stored alongside \code{NIG_ALPHA_0} and \code{NIG_BETA_0} in \code{constants.py}.

\paragraph{Change 3: Planning-scan termination criterion.}

The current \code{precompute_plan} loop calls \code{get_state_space}, which defines the upper scanning bound from the 99.9th percentile of a Normal distribution fit to the current belief. Under Stage~B this Gaussian heuristic must not be carried over unchanged. In the near-Cauchy regime ($\alpha_0^{\mathrm{NIG}} \approx 1$, small $t$), the Student-$t$ 99.9th percentile can be an order of magnitude larger than the distribution's scale, producing an excessively slow scan or an ill-bounded upper limit. Replace the distribution-based criterion with a fixed clinically safe overlap bound combined with the existing policy-flat termination condition (policy fallen to the minimum deliverable dose). This decouples scan termination from the choice of predictive distribution and does not require re-validation when the predictive changes.

\clearpage
\section{Required Numerical Design}

\subsection{State Discretisation}

Stage~B uses the same $(\mu, \sigma)$ grids and the same \code{_VOLUME_SPACE} as Stage~A (Section~3.2). The belief-state resolution, value-function shape, and peak memory footprint are unchanged.

\subsection{Complexity}

Compared to Stage~A, Stage~B introduces:
\begin{itemize}
\item A fraction-indexed branch-probability table of $\approx 29$\,MB at module load (versus $\approx 6$\,MB for the Stage~A table). This is a one-time cost per process invocation.
\item Student-$t$ CDF evaluation in place of Normal CDF during the precomputation step. \code{scipy.stats.t.cdf} is approximately $3$--$5\times$ slower per call than \code{scipy.stats.norm.cdf}; the precomputation time grows from near-instant to a few seconds, but this cost is not incurred per patient.
\item Per-node NIG parameter computation (five arithmetic operations per $(m_i, s_i, t)$ triple) when calling \code{_bellman_expectation} with the planning-time Student-$t$ distribution. This is negligible.
\end{itemize}

Per-patient DP runtime and peak memory are expected to remain within the Stage~A budget ($\approx35$\,s and $\approx2.75$\,GB), because the per-patient Bellman loop is structurally unchanged and accesses an equally-sized branch-probability slice.

\clearpage
\section{Validation and Acceptance Criteria}

\subsection{Implementation Correctness Gates}
\begin{enumerate}
\item Mathematical consistency: NIG posterior updates, predictive distribution, and branch-probability tables in code match this document.
\item Functional correctness: full test suite passes, and new tests cover fraction-dependent branch probabilities, NIG initialisation, and the equivalence of the Welford update with the NIG sufficient-statistics update.
\item Baseline lock-in: Stage~A benchmark artifacts (58-patient cohort replay) are frozen before Stage~B development begins and serve as the exclusive reference for all Stage~B comparisons.
\item Reproducibility: deterministic outputs under fixed discretisation and random settings.
\end{enumerate}

\subsection{Evaluation Metrics}

For each trajectory $i$ on the fixed 58-patient replay cohort, define:
\begin{itemize}
\item $B_i^{\mathrm{A}}$: Stage~A benefit,
\item $B_i^{\mathrm{B}}$: Stage~B benefit,
\item $B_i^{\mathrm{ideal}}$: best-case upper-bound benefit.
\end{itemize}
All three are measured on the same trajectory under identical simulation conditions.

Primary per-trajectory improvement:
\begin{equation}
\Delta_i = B_i^{\mathrm{B}} - B_i^{\mathrm{A}}.
\end{equation}

Distance-to-ideal:
\begin{equation}
G_i^{\mathrm{A}} = B_i^{\mathrm{ideal}} - B_i^{\mathrm{A}},\qquad
G_i^{\mathrm{B}} = B_i^{\mathrm{ideal}} - B_i^{\mathrm{B}}.
\end{equation}

Normalised gap closure:
\begin{equation}
R_i = \frac{B_i^{\mathrm{B}} - B_i^{\mathrm{A}}}{\max\!\bigl(B_i^{\mathrm{ideal}} - B_i^{\mathrm{A}},\;\epsilon\bigr)},
\end{equation}
with small $\epsilon > 0$ for numerical safety. $R_i = 0$ means no progress over Stage~A; $R_i = 1$ means the full Stage~A-to-ideal gap is closed.

\subsection{Acceptance Criteria}
\begin{enumerate}
\item Mean improvement is positive: $\mathbb{E}[\Delta_i] > 0$ with a positive paired-bootstrap 95\% confidence-interval lower bound.
\item Gap-to-ideal shrinks on average: $\mathbb{E}[G_i^{\mathrm{B}}] < \mathbb{E}[G_i^{\mathrm{A}}]$ and median $R_i > 0$.
\item Tail safety is not degraded: the fraction of cases with $\Delta_i < -1$\,ccGy does not increase relative to Stage~A.
\item Runtime and memory remain within the Stage~A budget.
\end{enumerate}

\section{Stage B Evaluation Protocol}

Stage~B must be evaluated in the following order before any model-selection or simplification decisions are made:

\begin{enumerate}
\item Run the full 58-patient replay benchmark under both Stage~A and Stage~B from the same frozen artifacts.
\item If the mean improvement is negligible ($|\mathbb{E}[\Delta_i]| < 0.01$\,ccGy): record this as a validated negative result. Stage~A is accepted as the production model and Stage~B is not promoted.
\item If the improvement is meaningful: compare in detail across:
\begin{enumerate}
\item expected total penalty and mean $\Delta_i$,
\item tail penalty at p90 and p95,
\item per-fraction dose trajectory differences,
\item terminal feasibility (underdose and overdose incidence),
\item runtime and peak memory.
\end{enumerate}
\item Validate generalisation via leave-one-out cross-validation: for each patient, refit the NIG hyperparameters on the remaining 57 patients and evaluate on the held-out patient. Compare the LOO $\mathbb{E}[\Delta_i]$ with the in-sample estimate from step~1. If the two are within noise, proceed with the full-cohort fit for the final model. If they diverge, investigate whether the hyperparameters are absorbing cohort-specific structure before promoting Stage~B. This step applies regardless of whether the improvement in steps~2--3 is negligible or meaningful. Leave-one-out is straightforward to implement because the marginal likelihood factorises over patients and each patient's contribution can be excluded by omitting it from the sum.
\end{enumerate}

Any subsequent simplification of Stage~B — for example reverting to a Normal predictive for fractions $t \geq t^*$ where the Student-$t$ and Normal are empirically indistinguishable — must pass the same comparison protocol against the full Stage~B result before being accepted.

\clearpage
\section{Stage B: Empirical Record (Log-NIG Implementation, March 2026)}

This section records what was implemented, what failed, what was fixed, and what was
ultimately concluded during the first concrete Stage~B development effort
(\code{feat/belief-state-dp} branch, March 2026).

\subsection{Model Choice: Log-Space NIG}

The original Stage~B plan (Section~3) assumed the NIG would operate on raw
overlap values in cc.  When initial fitting of the NIG hyperparameters was
attempted, the unconstrained optimum for $\alpha_0$ fell below the required floor
of~1.0, indicating that the symmetric Gaussian assumption of the NIG conflicts with
the right-skewed overlap distribution ($\text{skewness} = 2.73$ on the 58-patient
ACTION cohort).  The NIG model was therefore fitted in log-space: each observation
is transformed as $y_i = \log(v_i + \varepsilon)$ with $\varepsilon = 0.1$\,cc
before the NIG update.  In log-space the distribution is near-symmetric
($\text{skewness} = -0.15$), the NIG unimodal assumption is valid, and the
unconstrained optimum for $\alpha_0$ is interior ($\hat\alpha_0 = 2.74$, giving
degrees of freedom $2\alpha_0 + 1 = 6.49$ at fraction~1).

The log-NIG model was fitted by empirical Bayes: maximise the sequential NIG
marginal log-likelihood
\[
  \sum_{\text{patients}} \sum_{t=1}^{n} \log p\!\left(y_t \mid y_1,\dots,y_{t-1}\right)
\]
over the 58-patient cohort using L-BFGS-B with the reparameterisation described in
\code{scripts/fit\_nig\_hyperparams.py} (\code{-{}-log-space} flag).  The fitted
hyperparameters are stored in \code{constants.py} as \code{NIG\_LOG\_*}:

\begin{center}
\small
\begin{tabular}{lll}
\hline
\textbf{Constant} & \textbf{Value} & \textbf{Notes} \\
\hline
\code{NIG\_LOG\_OFFSET}  & $0.1$\,cc   & $\varepsilon$ in $\log(v+\varepsilon)$ \\
\code{NIG\_LOG\_MU\_0}   & $0.5230$\,nat  & Prior log-mean \\
\code{NIG\_LOG\_KAPPA\_0}& $0.1065$    & Prior pseudo-count on mean \\
\code{NIG\_LOG\_ALPHA\_0}& $2.7438$    & Interior solution; $\text{dof}_1 = 6.49$ \\
\code{NIG\_LOG\_BETA\_0} & $0.5526$\,nat$^2$ & Prior log-scale \\
\hline
\end{tabular}
\normalsize
\end{center}

The belief-state grid for log-NIG uses log-space coordinates: $\bar s_t = \text{running
mean of } y_i$ (nats) and $\sigma^{\log}_t = \sqrt{M_2/t}$ (running log-standard
deviation, nats).  The grid covers $[-2.5, 3.5]$\,nat in 100 non-uniform $\mu$
points and $[0.05, 1.50]$\,nat in 10 $\sigma$ points
(\code{\_MU\_GRID\_NIG}, \code{\_SIGMA\_GRID\_NIG} in \code{belief\_model.py}).

\subsection{Bug: Incorrect Grid Coordinates in \texttt{\_resolve\_current\_fraction}}

The initial log-NIG benchmark produced results worse than Stage~A by
$-0.235$\,ccGy (mean, 58-patient cohort).  Investigation revealed a grid-indexing
bug in \code{core\_adaptfx.py}: \code{\_resolve\_current\_fraction} was passing
the NIG \emph{posterior parameters} $(\mu_n, \hat\sigma_{\mathrm{MAP}})$ as the
DP grid lookup coordinates, whereas the precomputed branch-probability table
\code{\_P\_BELIEF\_NIG\_CACHE} is indexed by the \emph{Welford sufficient statistics}
$(\bar s_{\log}, \sigma^{\log}_{\mathrm{running}})$.  At fraction~2 with two
log-observations, this mismatch produced approximately a $2\times$ error in the
$\sigma$ coordinate.  The fix was to pass
$(\bar s_{\log},\, \max(\sqrt{M_2/n},\, \sigma^{\log}_{\min}))$
directly as the grid coordinates.  After the fix, the log-NIG benchmark was
$-0.147$\,ccGy vs.\ master and $-0.235$\,ccGy vs.\ Stage~A --- still worse than
Stage~A.

\subsection{Root Cause: NIG Cross-Term Inflates Predictive Scale}

With the correct grid coordinates, Stage~B remained worse than Stage~A.
Diagnosis identified the NIG cross-term in $\beta_n$:
\[
  \beta_n = \beta_0 + \tfrac{M_2}{2}
  + \underbrace{\frac{\kappa_0\,n\,(\bar s - \mu_0)^2}{2\kappa_n}}_{\text{cross-term}},
\]
as the structural cause.  For any patient whose log-mean $\bar s$ deviates from the
prior mean $\mu_0 = 0.523$\,nat (corresponding to $\approx 1.6$\,cc), this term
inflates $\beta_n$ and therefore the Student-$t$ predictive scale.  Patients with
near-zero overlaps --- which are far below the population mean --- receive a
substantially wider predictive than their data warrants, causing the DP to
over-hedge: it allocates doses conservatively even when it should exploit known-low
overlap fractions.

Stage~A avoids this because its Gamma prior on $\sigma$ is independent of the mean
estimate and is not inflated by deviations from a population prior.

\subsection{Fix Attempts and Results}

Three approaches were tried to remove or suppress the cross-term effect.

\paragraph{Option~0: Suppress the cross-term by setting $\kappa_0 \to 0$.}
With $\kappa_0 = 0.001$, the cross-term contributes at most
$0.001 \times n \times (\bar s - \mu_0)^2 / 2 \lesssim 0.003$\,nat$^2$ for any
patient --- negligible relative to $\beta_0 + M_2/2$.  This gives a closed-form
Student-$t$ predictive with scale $\approx \sqrt{2\beta_n / \alpha_n}$ that is
essentially free of shrinkage toward the population mean.  Result: mean delta vs.\
Stage~A of $+0.006$\,ccGy (2 patients improved, 2 worse, 54 tied out of 58).

\paragraph{Option~1: Re-fit $(\mu_0, \alpha_0, \beta_0)$ with $\kappa_0 = 0.001$
fixed (marginal-likelihood objective).}
Fitting with the cross-term suppressed changes the marginal-likelihood landscape;
the unconstrained optimum shifts to $\alpha_0 = 1.527$, $\beta_0 = 0.185$
(tighter prior scale, heavier tails).  Result: mean delta vs.\ Stage~A
$+0.006$\,ccGy with \emph{different} patients changing direction --- no net
improvement over Option~0.

\paragraph{Option~2: Grid search over $(\alpha_0, \beta_0)$ optimising DP benefit
directly.}
A two-dimensional grid search (10 $\alpha_0$ values $\times$ 8 prior-scale values
$= 80$ combinations, parameterised by $s_1 = \sqrt{2\beta_0/(\alpha_0+0.5)}$ to
cover the space uniformly) was run on all 58 patients with $\kappa_0 = 0.001$
fixed.  After evaluating 36 of 80 combinations (covering $\alpha_0 \in
[0.55, 1.527]$ fully), every combination produced a negative mean delta vs.\
Stage~A, with the best being $-0.030$\,ccGy at $(\alpha_0 = 0.75,\,
s_1 = 0.25\,\text{nat})$.  The search was stopped at this point.

\subsection{Conclusions}

\begin{enumerate}
\item The log-NIG model with $\kappa_0 = 0.001$ achieves approximate parity with
Stage~A ($+0.006$\,ccGy mean, within the grid-quantisation noise of the
$\pm0.005$\,ccGy floor established during the Stage~A $N_\mu$ sweep).  It does not
meaningfully exceed Stage~A.

\item Direct optimisation of DP benefit over hyperparameters did not find any
$(\alpha_0, \beta_0)$ combination that beats Stage~A.  The pattern across 36 grid
points is uniformly negative, with no sign of an unexplored positive region.

\item The 54/58 tied patients have dose trajectories that are identical under
Stage~A and log-NIG regardless of hyperparameters.  The tying mechanism is
\emph{not} that patients hit the hard min/max dose bounds; it is that for patients
with consistent overlap patterns, both the Gaussian (cc-space) and Student-$t$
(log-space) predictives lead to the same qualitative conclusion at every fraction ---
``overlap is reliably low, push dose high'' or ``overlap is reliably high, keep
dose low'' --- so the Bellman recursion selects the same 0.5\,Gy dose bin either
way.  A belief model can only change a dose decision when the predictive uncertainty
about future overlaps is large enough to shift the optimal action across a dose-step
boundary.  Improvement can therefore only come from the 4~patients whose overlap
trajectories fall in the ambiguous range where model choice actually crosses such a
boundary: the Student-$t$ tails help high-variability patients (e.g.\ patient~11,
$\hat\sigma \approx 0.26$\,cc) and hurt near-zero-overlap patients (e.g.\ patients
22 and~41, mean $< 0.15$\,cc) via the log-floor effect
($\log(0 + 0.1) \approx -2.30$\,nat), netting close to zero.  Any belief model
improvement will be bounded by the number of patients in this ambiguous range unless
the model change is structurally different enough to move additional patients off the
tying point.

\item \textbf{The semi-conjugate model (independent Normal prior on $\mu$,
Inverse-Gamma prior on $\sigma^2$) is mathematically equivalent to $\kappa_0 \to
0$ for this predictive family.}  Since $\kappa_0 = 0.001$ already approximates this
limit to within $0.003$\,nat$^2$ in $\beta_n$, implementing the semi-conjugate
model in full generality would reproduce the Option~0 result.  It is not a
promising avenue.

\item The fundamental obstacle is that the log-Student-$t$ predictive in
log-space does not systematically improve DP dose decisions over the Gaussian
predictive in cc-space for this cohort.  The log-transform floor
($v = 0 \mapsto \log 0.1 \approx -2.30$\,nat$\approx -2.30$\,nat) clusters near-zero
overlap observations artificially, making Stage~A's cc-space Gaussian more
naturally concentrated near zero for the near-zero-overlap patients.
\end{enumerate}

\subsection{Status of Stage B in the Codebase}

The log-NIG implementation (\code{use\_nig=True}) is present in
\code{core\_adaptfx.py} and \code{belief\_model.py} and passes all non-slow tests.
It is \emph{not} the default code path (\code{use\_nig=False} by default).  The
implementation is correct (bug fixed, grid coordinates consistent) and the
hyperparameters in \code{constants.py} reflect the best configuration found
($\kappa_0 = 0.001$, empirically Bayes-fitted $\alpha_0, \beta_0$).  Given the
evaluation results above, Stage~B in its current form is not recommended for
production use.  Any future attempt to exceed Stage~A should consider:

\begin{itemize}
\item A different transform or mixture model that avoids the log-floor problem for
near-zero-overlap patients while retaining the right-skew correction for typical
patients.
\item Directly optimising the predictive distribution for DP performance rather than
marginal likelihood, with early stopping guided by the 58-patient benchmark rather
than by convergence of the likelihood objective.
\item A fundamentally different belief model (e.g.\ non-parametric or copula-based)
that does not assume exchangeability across fractions.
\end{itemize}

\end{document}
